\subsection{Witt decomposition.}

In addition to the conventions already in force since the beginning of the present paragraph, we shall add the following:

Condition (T). --- For all \(x \in E\) , there exists \(\alpha \in A\) such that \(\Phi(x, x) = \alpha + \varepsilon\bar{\alpha}\) .

This condition is always satisfied when \(\Phi\) is alternating, or when \(\varepsilon = 1\) and A is a field of characteristic \(\neq 2\) , taking then \(\alpha = \frac{1}{2}\Phi(x, x)\) (cf. Exercises 1 and 14).

Lemma 1. --- Let \(\Phi\) be a form \(\varepsilon\) -hermitian satisfying (T) (resp. the bilinear form associated with a quadratic form Q) on E, and let F be a totally isotropic (resp. totally singular) subspace of E, not reduced to 0. For every \(x \in E\) not orthogonal to F and all \(\alpha \in A\) , there exists \(y \in F\) such that

\[
\Phi (x + y, x + y) = \alpha + \varepsilon \bar {\alpha} \quad (\text { resp. } Q (x + y) = \alpha).
\]

Indeed, let us set \(\Phi(x, x) = \beta + \varepsilon\bar{\beta}\) (resp. \(Q(x) = \beta\) ). For \(y \in F\) we then have \(\Phi(x + y, x + y) = (\beta + \Phi(x, y)) + \varepsilon(\overline{\beta + \Phi(x, y)})\) since \(\Phi(y, y) = 0\) (resp. \(Q(x + y) = \beta + \Phi(x, y)\) since \(Q(y) = 0\) ). Since \(x\) is not orthogonal to \(F\) , the affine linear function \(y \to \Phi(x, y) + \beta\) on \(F\) is not constant; it therefore takes the value \(\alpha\) for some element \(y\) of \(F\) , which thus answers the question.

A Witt decomposition of E is any decomposition of E as a direct sum of three subspaces F, F', G such that F and F' are totally isotropic (resp. totally singular) and G is non-isotropic and orthogonal to F + F'; if E is finite-dimensional, the matrix of Φ with respect to a basis of E adapted to a Witt decomposition of E takes the form

\[
\left( \begin{array}{c c c} 0 & U & 0 \\ \varepsilon \overline {{U}} & 0 & 0 \\ 0 & 0 & V \end{array} \right)\tag{1}
\]

We say that \(\Phi\) is a neutral form if it is nondegenerate and if E is finite-dimensional and is the direct sum of two totally isotropic (resp. totally singular) subspaces. The direct sum of two neutral forms is a neutral form.

PROPOSITION 2. --- Let Φ be a nondegenerate ε-hermitian form satisfying (T) (resp. the bilinear form associated with a nondegenerate quadratic form Q), and let F be a totally isotropic (resp. totally singular) subspace of finite dimension r.

\begin{enumerate}
\def\labelenumi{\alph{enumi})}
\item
  If \(F'\) is a totally isotropic subspace of dimension r such that \(F' \cap F^{0} = \{0\}\) , then \(F + F'\) is non-isotropic and, for any basis \((f_{i})\) of F, there exists a basis \((f_{i}')\) of \(F'\) such that \(\Phi(f_{i}, f_{j}') = \delta_{ij}\) (Kronecker index) for \(i, j = 1, \ldots, r\) .
\item
  If G is a totally isotropic (resp. totally singular) subspace of dimension \(\leqslant r\) such that \(G \cap F^{0} = \{0\}\) , there exists a totally isotropic (resp. totally singular) subspace \(F' \supset G\) of dimension r such that \(F' \cap F^{0} = \{0\}\) .
\end{enumerate}

Let \(\Psi\) be the restriction of \(\Phi\) to \(F \times F'\) ; for \(x' \in F'\) , the relation \(\langle \Phi(x, x') = 0\) for all \(x \in F\rangle\) implies \(x = 0\) since \(F' \cap F^0 = \{0\}\) . The assertion \(a)\) then follows from the corollary of Proposition 6 of §1, no. 6, except for the fact that \(F + F'\) is non-isotropic. Now the subspace \(H = (F + F') \cap (F + F')^0\) is equal to \((F + F') \cap F^0 \cap F'^0\) . Since \(F \subset F^0\) , one has \((F + F') \cap F^0 = F + (F' \cap F^0) = F\) , whence \(H = F \cap F'^0\) ; hence \(H = \{0\}\) since we have seen that \(\Psi\) is nondegenerate. This indeed proves that \(F + F'\) is non-isotropic.

To demonstrate \(b\) ), we proceed by descending induction on \(s = \dim G\) . It thus suffices for us to prove that, if \(s < r\) , there exists a totally isotropic (resp. totally singular) subspace \(G'\) containing \(G\) , of dimension \(s + 1\) , and such that \(G' \cap F^0 = \{0\}\) . Since \(\dim G < \dim F\) , the restriction of \(\Phi\) to \(F \times G\) is degenerate, and since \(G \cap F^0\) is zero, \(F \cap G^0\) is nonzero. If one then had \(G + F^0 \supset G^0\) one would deduce, by taking the orthogonal subspaces and noting that \(F = F^{00}\) and \(G = G^{00}\) ( \(\S 1\) , no. 6, cor. 1 of prop. 4), that \(G^0 \cap F \subset G\) , whence

\[
\mathrm{G} ^ {0} \cap \mathrm{F} \subset \mathrm{G} \cap \mathrm{F} \subset \mathrm{G} \cap \mathrm{F} ^ {0} = \{0 \},
\]

which is impossible. There then exists an element \(x\) of \(G^0\) such that \(x \notin G + F^0\) ; since \(F \subset F^0\) , one can add to \(x\) a vector of \(G^0 \cap F\) without modifying these properties; since \(G^0 \cap F\) is totally isotropic (resp. totally singular) and \(\neq \{0\}\) , lemma 1 shows that one can choose \(x\) isotropic (resp. singular).

Then the subspace \(G' = G + Ax\) is of dimension \(s + 1\) , and is totally isotropic (resp. totally singular); moreover we have \(G' \cap F^{0} = \{0\}\) because, if \(y = z + ax\) ( \(z \in G, a \in A\) ) is in \(F^{0}\) , we have a = 0, since otherwise \(x \in F^{0} + G\) contrary to the choice of x, whence \(y = z \in G \cap F^{0} = \{0\}\) and y = 0. Consequently the subspace \(G'\) answers the question.

COROLLARY 1. --- If F is a totally isotropic (resp. totally singular) subspace of dimension r, there exists a totally isotropic (resp. totally singular) subspace F' of dimension r such that F ∩ F' = \{0\} and such that F + F' is non-isotropic.

It suffices to take G = \{0\} in prop. 2, b).

COROLLARY 2. --- Two neutral ε-hermitian forms on spaces of the same dimension over A are equivalent.

Remark. --- Under the conditions of corollary 1, E is the direct sum of F + F' and of the orthogonal of F + F'. One thus has a Witt decomposition of E. By prop. 2, a), there exist bases of F and F' such that, in the matrix (1) of Φ, the block U is the identity matrix 1.

PROPOSITION 3. --- Let \(\Phi\) be a nondegenerate \(\varepsilon\) -hermitian form satisfying (T) (resp. the bilinear form associated with a nondegenerate quadratic form Q). Let \(F_{1}\) and \(F_{2}\) be two maximal totally isotropic (resp. totally singular) subspaces of E, one of the two being finite-dimensional. Set \(F = F_{1} \cap F_{2}\) . Let \(S_{i}\) ( \(i = 1, 2\) ) be a supplement of F in \(F_{i}\) ; set \(S = S_{1} + S_{2}\) . There then exist two subspaces G and H of E such that

\begin{enumerate}
\def\labelenumi{\alph{enumi})}
\item
  The subspaces G + F, S, and H are non-isotropic and pairwise orthogonal;
\item
  E is the direct sum of F, S, G, and H;
\item
  There is no nonzero isotropic (resp. singular) vector in H;
\item
  G is totally isotropic (resp. totally singular).
\end{enumerate}

Additionally, \(F_{1}\) and \(F_{2}\) are both finite-dimensional and we have \(\dim \mathrm{F}_1 = \dim \mathrm{F}_2, \dim \mathrm{G} = \dim \mathrm{F}, \dim \mathrm{S}_1 = \dim \mathrm{S}_2, \operatorname{codim} \mathrm{H} = 2 \dim \mathrm{F}_1.\)

First, note that if N is a maximal totally isotropic (resp. totally singular) subspace, then every isotropic (resp. singular) vector x orthogonal to N belongs to N, for otherwise \(N + Ax\) would contradict the maximality of N. Thus, if, for i = 1, 2, \(x_{i}\) is an isotropic (resp. singular) vector of \(F_{i}^{0}\) , one has \(x_{i} \in F_{i}\) On the other hand, if y is an element of \(S_{1}\) orthogonal to \(S_{2}\) it is orthogonal to \(F_{1}\) since \(F_{1}\) is totally isotropic, hence to F, and consequently to \(F_{2} = S_{2} + F\) . Since y is isotropic (resp. singular) and is orthogonal to \(F_{2}\) , one has

\[
y \in S _ {1} \cap F _ {2} = S _ {1} \cap F _ {1} \cap F _ {2} = S _ {1} \cap F = \{0 \}.
\]

We therefore have \(S_1 \cap S_2^0 = \{0\}\) , and likewise \(S_2 \cap S_1^0 = \{0\}\) . Since one of the two subspaces \(F_1, F_2\) , for example \(F_1\) , is finite-dimensional, \(S_1\) is finite-dimensional, hence \(S_1^0\) has finite codimension (§ 1, no. 6, cor. 1 of prop. 4), and consequently \(S_2\) is finite-dimensional since \(S_2 \cap S_1^0 = \{0\}\) ; moreover this shows that one has dim \(S_2 \leqslant \text{codim } S_1^0 = \text{dim } S_1\) ; likewise dim \(S_1 \leqslant \text{dim } S_2\) , whence dim \(S_1 = \text{dim } S_2\) . Prop. 2 a) then shows that \(S = S_1 + S_2\) is non-isotropic.

This being so, the orthogonal N of S is non-isotropic (no. 1, cor. of prop. 1) and contains F; cor. 1 of prop. 2 therefore shows that there exists a subspace G totally isotropic (resp. totally singular) of N such that dim G = dim F, that G ∩ F = \{0\} and that G + F is non-isotropic. Thus d) is satisfied by G. One then satisfies a) and b) by taking for H the orthogonal of G + F in N. As for c), one notes that, since H is orthogonal to F₁ = S₁ + F, there is no nonzero isotropic (resp. singular) vector in H by virtue of what was remarked at the beginning of the proof and of the fact that H ∩ F₁ = \{0\}. Finally, certain of the assertions concerning dimensions were proved along the way; the others follow from them trivially.

COROLLARY 1. --- Under the hypotheses of prop. 3, two maximal totally isotropic (resp. totally singular) subspaces of finite dimension have the same dimension. For every maximal totally isotropic (resp. totally singular) subspace F of finite dimension, there exists another one \(\mathbf{F}'\) such that \(\mathbf{F} \cap \mathbf{F}' = \{0\}\) , and under these conditions \(\mathbf{F} + \mathbf{F}'\) is non-isotropic.

If \(F \cap F' = \{0\}\) , one will have \(G = \{0\}\) with the notations of prop. 3, and \(F + F'\) will be non-isotropic. The other assertions follow trivially from prop. 3 and cor. 1 of prop. 2.

COROLLARY 2. --- Let Q be a nondegenerate quadratic form on a vector space E of finite dimension n over an algebraically closed field A; there then exists a basis \((e_i)_{1 \leqslant i \leqslant n}\) of E such that

\begin{enumerate}
\def\labelenumi{(\arabic{enumi})}
\setcounter{enumi}{1}
\tightlist
\item
\end{enumerate}

\[
\mathrm{Q} (\sum_ {i = 1} ^ {n} x _ {i} e _ {i}) = \sum_ {i = 1} ^ {\nu} x _ {i} x _ {i + \nu} \quad \text {   if   } n = 2 \nu ,\tag{3}
\]

\[
\mathrm{Q} \left(\sum_ {i = 1} ^ {n} x _ {i} e _ {i}\right) = \sum_ {i = 1} ^ {\nu} x _ {i} x _ {i + \nu} + x _ {2 \nu + 1} ^ {2} \quad \text {   if   } n = 2 \nu + 1.
\]

Indeed, let \(F_{1}\) and \(F_{2}\) be two maximal totally singular subspaces such that \(F_{1} \cap F_{2} = \{0\}\) (cor. 1) and let q be their dimension. One then has \(G = \{0\}\) with the notations of prop. 3. Taking a basis \((e_{i})_{1 \leqslant i \leqslant q}\) of \(F_{1}\) and a basis \((e_{i})_{q+1 \leqslant i \leqslant 2q}\) of \(F_{2}\) such that \(\Phi(e_{i}, e_{j+q}) = \delta_{ij}\) for \(i, j = 1, \ldots, q\) (prop. 2 a)), we see that it suffices to show that dim \(H \leqslant 1\) . Now, if \(x \in H\) , \(y \in H\) and if \(x \neq 0\) , the equation \(Q(y - ax) = Q(y) - a\Phi(x, y) + a^{2}Q(x) = 0\) has at least one solution \(a_{0}\) since \(Q(x) \neq 0\) and one has \(y = a_{0}x\) since every singular vector of H is zero.

DEFINITION 3. --- Suppose that E is finite-dimensional and that \(\Phi\) is a nondegenerate \(\varepsilon\) -hermitian form satisfying (T) (resp. the bilinear form associated with a nondegenerate quadratic form Q). The index of \(\Phi\) (resp. Q) is the common dimension of the maximal totally isotropic (resp. totally singular) subspaces of E.

If n is the dimension of E and v the index of \(\Phi\) (resp. Q), Prop. 3 shows that one has

\[
n \geqslant 2 v.\tag{4}
\]

Moreover, since every totally isotropic (resp. totally singular) subspace is contained in a maximal totally isotropic (resp. totally singular) subspace, the maximal totally isotropic (resp. totally singular) subspaces are precisely those of dimension v. The assertion that \(\Phi\) (resp. Q) has index 0 means that every isotropic (resp. singular) vector of E is zero. In a space of even dimension \(n\) the neutral forms are those of index \(\frac{1}{2} n\) ; there is no neutral form in a space of odd dimension. Prop. 3 shows that every form is a direct sum of a neutral form and a form of index 0.

PROPOSITION 4. --- Let Q be a nondegenerate quadratic form on E such that there exists a vector \(x \neq 0\) of E with Q(x) = 0. For every element a of A, there then exists \(y \in E\) such that Q(y) = a.

Indeed, by cor. 1 of prop. 2, there exists a subspace \(\mathbf{G} = \mathbf{F} + \mathbf{F}'\) ( \(\mathbf{F}, \mathbf{F}'\) : totally singular subspaces of dimension 1) of \(\mathbf{E}\) of dimension 2, such that the restriction of \(\mathbf{Q}\) to \(\mathbf{G}\) is neutral. If \(\{e, e'\}\) ( \(e \in \mathbf{F}, e' \in \mathbf{F}'\) ) is a basis of \(\mathbf{G}\) , one has

\[
\mathrm{Q} (x e + x ^ {\prime} e ^ {\prime}) = b x x ^ {\prime} \quad (x \in \mathrm{A}, x ^ {\prime} \in \mathrm{A}, b \in \mathrm{A}, b \neq 0).
\]

It therefore suffices to take for y the vector \(ae + b^{-1}e'\) .
% semantic-fragments: legacy-input-seam
\relax{}
